\subsection{Abstand eines Punktes von einer Ebene}
\label{subsec:analytic_geometry_point_to_plane_distance}

Bei einer Ebene zeigt der Normalenvektor bereits in Richtung der kürzesten
Verbindung. Deshalb genügt eine skalare Projektion auf diese Richtung.

\begin{center}
  \CachedTikzFigure{content/resources/chapters/analytic_geometry/tikz/point_plane_distance.tex}
\end{center}

Der Abstand ist der Betrag der parallelen Komponente $\overrightarrow{PQ}_{\parallel}$,
also der orthogonale Projektion von $\overrightarrow{PQ}$ auf die
Normalenrichtung $\vec n$.
\begin{ClozeArea}[34]{3}
  \[
    d=\left\lvert\textcolor{black!75}{\overrightarrow{PQ}_{\parallel}}\right\rvert
    =\left\lvert
      \bigl(\textcolor{red!75!black}{\overrightarrow{PQ}}
      \cdot\textcolor{THMBlue}{\vec e_n}\bigr)\textcolor{THMBlue}{\vec e_n}
    \right\rvert
    =\left\lvert
      \textcolor{red!75!black}{\overrightarrow{PQ}}
      \cdot\textcolor{THMBlue}{\vec e_n}
    \right\rvert
      \overbrace{\left\lvert\textcolor{THMBlue}{\vec e_n}\right\rvert}^{=1}
    =\left\lvert
      \textcolor{red!75!black}{\overrightarrow{PQ}}
      \cdot\textcolor{THMBlue}{\vec e_n}
    \right\rvert.
  \]
\end{ClozeArea}

\begin{TheoremBox}{Punkt--Ebenen-Abstand}{analytic_geometry_point_plane_distance}
  Sei $P$ ein Punkt der Ebene $E$ und $Q$ ein beliebiger Punkt, 
  so ist der Abstand $d$ von $Q$ zur Ebene $E$ gegeben durch
  \[
    d
    =\left\lvert\overrightarrow{PQ}_{\parallel}\right\rvert
    =\left\lvert\overrightarrow{PQ}\cdot\vec e_n\right\rvert.
  \]
\end{TheoremBox}

\begin{ExampleBox}{Abstand eines Punktes von einer Ebene}{analytic_geometry_point_plane_distance_example}
  Gegeben sind die Ebene $E$ in Parameterform und der Punkt $Q(3,2,2)$:
  \[
    E:\quad
    \vec x=\overrightarrow{OP}+\lambda\vec a+\mu\vec b,
    \qquad
    \overrightarrow{OP}=\begin{pmatrix}1\\0\\0\end{pmatrix},
    \quad
    \vec a=\begin{pmatrix}0\\2\\-1\end{pmatrix},
    \quad
    \vec b=\begin{pmatrix}0\\1\\2\end{pmatrix}.
  \]
  Aus den Richtungsvektoren bestimmen wir einen
  Einheitsnormalenvektor:
  \begin{ClozeArea}[34]{11}
    \[
      \vec n
      =\vec a\times\vec b
      =\begin{pmatrix}0\\2\\-1\end{pmatrix}
        \times\begin{pmatrix}0\\1\\2\end{pmatrix}
      =\begin{pmatrix}5\\0\\0\end{pmatrix}
      \quad\Longrightarrow\quad
      \vec e_n=\begin{pmatrix}1\\0\\0\end{pmatrix}.
    \]
  \end{ClozeArea}
  Wir setzen in die Formel aus Satz~\ref{thm:analytic_geometry_point_plane_distance} ein:
  \begin{ClozeArea}[34]{4}
    \[
      d
      =\left\lvert\overrightarrow{PQ}\cdot\vec e_n\right\rvert
      =\left\lvert
        \bigl(\overrightarrow{OQ}-\overrightarrow{OP}\bigr)\cdot\vec e_n
      \right\rvert
      =\left\lvert
        \left(
          \begin{pmatrix}3\\2\\2\end{pmatrix}
          -\begin{pmatrix}1\\0\\0\end{pmatrix}
        \right)\cdot
        \begin{pmatrix}1\\0\\0\end{pmatrix}
      \right\rvert
      =2.
    \]
  \end{ClozeArea}
  Der Punkt $Q$ hat somit den Abstand $\Cloze{3cm}{2}$ von der Ebene.
\end{ExampleBox}
